\chapter{Whiteboard Exercises and Worked Examples}
\label{ch:whiteboard_exercises}

Technical interviews frequently require solving concrete mathematical or algorithmic problems on a whiteboard or virtual notepad. This chapter provides eight end-to-end worked problems directly reflecting the mathematics, algorithms, and engineering logic of BengaluruTrafficAI. Each exercise states explicit assumptions, step-by-step arithmetic derivations, relevant code references, and edge-case evaluations.

\begin{alertnote}
\textbf{Note on Example Data:} Numerical inputs and coordinates in these exercises are formulated as illustrative examples to demonstrate algorithmic mechanics clearly.
\end{alertnote}

% -----------------------------------------------------------------------------
\section{Exercise 1: Computing Intersection over Union (IoU)}
\label{sec:wb_iou}
% -----------------------------------------------------------------------------

\begin{interviewq}{Whiteboard Problem 1: Two-Box IoU Derivation}
\textbf{Problem Statement:}
In \texttt{ai/tracker\_manager.py}, vehicle tracking associates incoming YOLO detections with existing tracks by calculating pairwise bounding box Intersection over Union (IoU). Given:
\begin{itemize}
    \item Existing Track Box $B_1 = [x_1, y_1, x_2, y_2] = [100, 150, 300, 450]$
    \item New Detection Box $B_2 = [x_1, y_1, x_2, y_2] = [160, 200, 360, 500]$
\end{itemize}
Derive the intersection area, union area, and IoU value. State whether these boxes associate under our tracking threshold of $\theta_{\text{IoU}} = 0.3$.
\tcblower
\textbf{Step 1: Compute Individual Box Areas}
\begin{align*}
W_1 &= x_2 - x_1 = 300 - 100 = 200 \text{ px} \\
H_1 &= y_2 - y_1 = 450 - 150 = 300 \text{ px} \\
\text{Area}(B_1) &= W_1 \times H_1 = 200 \times 300 = 60{,}000 \text{ px}^2 \\[6pt]
W_2 &= x_2 - x_1 = 360 - 160 = 200 \text{ px} \\
H_2 &= y_2 - y_1 = 500 - 200 = 300 \text{ px} \\
\text{Area}(B_2) &= W_2 \times H_2 = 200 \times 300 = 60{,}000 \text{ px}^2
\end{align*}

\textbf{Step 2: Compute Intersection Bounding Coordinates}
The intersection box $B_{\cap} = [x_{\cap 1}, y_{\cap 1}, x_{\cap 2}, y_{\cap 2}]$ is given by:
\begin{align*}
x_{\cap 1} &= \max(x_{1,1}, x_{2,1}) = \max(100, 160) = 160 \\
y_{\cap 1} &= \max(y_{1,1}, y_{2,1}) = \max(150, 200) = 200 \\
x_{\cap 2} &= \min(x_{1,2}, x_{2,2}) = \min(300, 360) = 300 \\
y_{\cap 2} &= \min(y_{1,2}, y_{2,2}) = \min(450, 500) = 450
\end{align*}

\textbf{Step 3: Compute Intersection Area}
\begin{align*}
W_{\cap} &= \max(0, x_{\cap 2} - x_{\cap 1}) = \max(0, 300 - 160) = 140 \text{ px} \\
H_{\cap} &= \max(0, y_{\cap 2} - y_{\cap 1}) = \max(0, 450 - 200) = 250 \text{ px} \\
\text{Area}(B_{\cap}) &= W_{\cap} \times H_{\cap} = 140 \times 250 = 35{,}000 \text{ px}^2
\end{align*}

\textbf{Step 4: Compute Union Area and IoU}
\begin{align*}
\text{Area}(B_{\cup}) &= \text{Area}(B_1) + \text{Area}(B_2) - \text{Area}(B_{\cap}) \\
&= 60{,}000 + 60{,}000 - 35{,}000 = 85{,}000 \text{ px}^2 \\[6pt]
\text{IoU}(B_1, B_2) &= \frac{\text{Area}(B_{\cap})}{\text{Area}(B_{\cup})} = \frac{35{,}000}{85{,}000} \approx \mathbf{0.4118}
\end{align*}

\textbf{Conclusion:} Since $\text{IoU} \approx 0.412 \ge 0.3$, the boxes successfully match. The tracker updates Track 1 with Detection 2's bounding box and resets its missing frame counter to 0.
\end{interviewq}

% -----------------------------------------------------------------------------
\section{Exercise 2: Greedy Bipartite Matching for Multi-Object Tracking}
\label{sec:wb_matching}
% -----------------------------------------------------------------------------

\begin{interviewq}{Whiteboard Problem 2: Tracking Association Matrix}
\textbf{Problem Statement:}
Assume at frame $t$, the system has two active tracks: $T_1$ and $T_2$. Frame $t$ produces three YOLO detections: $D_1, D_2, D_3$. The pairwise IoU cost matrix is computed as follows:
\begin{equation*}
\mathbf{C}_{\text{IoU}} =
\begin{pmatrix}
0.65 & 0.12 & 0.05 \\
0.18 & 0.72 & 0.28
\end{pmatrix}
\end{equation*}
where rows represent $[T_1, T_2]$ and columns represent $[D_1, D_2, D_3]$. Using the greedy matching algorithm implemented in \texttt{ai/tracker\_manager.py} with threshold $\theta = 0.3$, trace:
\begin{enumerate}
    \item Matched track-detection pairs.
    \item Unmatched detections and action taken.
    \item Unmatched tracks and action taken.
\end{enumerate}
\tcblower
\textbf{Step 1: Greedy Selection by Maximum Overlap}
\begin{enumerate}
    \item Global maximum in matrix is $\mathbf{0.72}$ at index $(T_2, D_2)$. Since $0.72 \ge 0.3$:
    \begin{itemize}
        \item Pair $(T_2, D_2)$ is matched.
        \item Remove row $T_2$ and column $D_2$ from remaining candidates.
    \end{itemize}
    \item Next maximum in remaining elements is $\mathbf{0.65}$ at index $(T_1, D_1)$. Since $0.65 \ge 0.3$:
    \begin{itemize}
        \item Pair $(T_1, D_1)$ is matched.
        \item Remove row $T_1$ and column $D_1$.
    \end{itemize}
    \item No active tracks remain to be matched.
\end{enumerate}

\textbf{Step 2: Process Unmatched Candidates}
\begin{itemize}
    \item \textbf{Matched Tracks:} $T_1 \leftarrow D_1$ and $T_2 \leftarrow D_2$. Their bounding boxes and positions update, and their \texttt{missing\_frames} count resets to 0.
    \item \textbf{Unmatched Detections:} $D_3$ had an overlap of at most 0.28 (with $T_2$, which was already claimed) and 0.05 with $T_1$. $D_3$ is identified as a new vehicle entering the intersection. The tracker allocates a new unique track ID (e.g., $T_3$) and registers its centroid.
    \item \textbf{Unmatched Tracks:} None in this frame. If any track remained unmatched, its \texttt{missing\_frames} counter would increment by 1.
\end{itemize}
\end{interviewq}

% -----------------------------------------------------------------------------
\section{Exercise 3: Calculating Speed from Pixel Displacement}
\label{sec:wb_speed}
% -----------------------------------------------------------------------------

\begin{interviewq}{Whiteboard Problem 3: Speed Estimation}
\textbf{Problem Statement:}
In \texttt{ai/speed\_estimator.py}, vehicle speed is estimated from tracked centroid motion. Assume:
\begin{itemize}
    \item Frame $t_1$: Timestamp $\tau_1 = 100.00$ s, vehicle centroid $C_1 = (320, 240)$ px.
    \item Frame $t_2$: Timestamp $\tau_2 = 100.40$ s, vehicle centroid $C_2 = (320, 310)$ px.
    \item Camera scaling constant: $k = 0.08$ meters/pixel.
    \item Physical speed clamping limits: $[0.0, 140.0]$ km/h.
\end{itemize}
Calculate the vehicle speed in m/s and km/h. Does it trigger an overspeed violation if the posted limit is $50$ km/h?
\tcblower
\textbf{Step 1: Calculate Pixel Displacement and Elapsed Time}
\begin{align*}
\Delta x &= x_2 - x_1 = 320 - 320 = 0 \text{ px} \\
\Delta y &= y_2 - y_1 = 310 - 240 = 70 \text{ px} \\
d_{\text{px}} &= \sqrt{(\Delta x)^2 + (\Delta y)^2} = \sqrt{0 + 70^2} = 70 \text{ px} \\[6pt]
\Delta \tau &= \tau_2 - \tau_1 = 100.40 - 100.00 = 0.40 \text{ s}
\end{align*}

\textbf{Step 2: Convert to Physical Distance and Velocity}
\begin{align*}
d_{\text{meters}} &= d_{\text{px}} \times k = 70 \times 0.08 = 5.60 \text{ meters} \\[6pt]
v_{\text{m/s}} &= \frac{d_{\text{meters}}}{\Delta \tau} = \frac{5.60}{0.40} = 14.0 \text{ m/s} \\[6pt]
v_{\text{km/h}} &= v_{\text{m/s}} \times 3.6 = 14.0 \times 3.6 = \mathbf{50.4 \text{ km/h}}
\end{align*}

\textbf{Step 3: Evaluation Against Bounds and Limits}
\begin{itemize}
    \item \textbf{Sanity Bound Check:} $0.0 \le 50.4 \le 140.0$ km/h $\implies$ valid measurement (not rejected as noise).
    \item \textbf{Speed Limit Evaluation:} $50.4 > 50.0$ km/h $\implies$ marginally exceeds speed limit threshold.
    \item \textbf{Interview Follow-Up / Caveat:} In a real interview, emphasize that a 0.4 km/h delta is within calibration error margin for uncalibrated perspective projection, illustrating the necessity of homography calibration.
\end{itemize}
\end{interviewq}

% -----------------------------------------------------------------------------
\section{Exercise 4: Allocating Dynamic Signal Green Time}
\label{sec:wb_signal}
% -----------------------------------------------------------------------------

\begin{interviewq}{Whiteboard Problem 4: Signal Green-Time Allocation}
\textbf{Problem Statement:}
In \texttt{backend/signal\_service.py}, traffic signal green times are allocated proportionally based on incoming lane vehicle counts. Consider a 4-way intersection (North, South, East, West) with:
\begin{itemize}
    \item Vehicle counts: $N_{\text{North}} = 45$, $N_{\text{South}} = 30$, $N_{\text{East}} = 15$, $N_{\text{West}} = 10$.
    \item Total target green cycle budget: $T_{\text{green}} = 120$ seconds.
    \item Minimum green guard: $G_{\min} = 10$ seconds.
    \item Maximum green clamp: $G_{\max} = 60$ seconds.
\end{itemize}
Derive the proportional green time for each approach and verify cycle boundary constraints.
\tcblower
\textbf{Step 1: Compute Total Vehicle Count}
\begin{equation*}
N_{\text{total}} = 45 + 30 + 15 + 10 = 100 \text{ vehicles}
\end{equation*}

\textbf{Step 2: Compute Raw Proportional Allocations}
\begin{align*}
G_{\text{North, raw}} &= 120 \times \left(\frac{45}{100}\right) = 120 \times 0.45 = \mathbf{54.0 \text{ s}} \\
G_{\text{South, raw}} &= 120 \times \left(\frac{30}{100}\right) = 120 \times 0.30 = \mathbf{36.0 \text{ s}} \\
G_{\text{East, raw}}  &= 120 \times \left(\frac{15}{100}\right) = 120 \times 0.15 = \mathbf{18.0 \text{ s}} \\
G_{\text{West, raw}}  &= 120 \times \left(\frac{10}{100}\right) = 120 \times 0.10 = \mathbf{12.0 \text{ s}}
\end{align*}

\textbf{Step 3: Apply Boundary Clamping}
Every raw allocation $G_i$ must satisfy $G_{\min} \le G_i \le G_{\max}$:
\begin{itemize}
    \item North: $10 \le 54.0 \le 60 \implies 54.0$ s (Valid)
    \item South: $10 \le 36.0 \le 60 \implies 36.0$ s (Valid)
    \item East: $10 \le 18.0 \le 60 \implies 18.0$ s (Valid)
    \item West: $10 \le 12.0 \le 60 \implies 12.0$ s (Valid)
\end{itemize}
Sum check: $54 + 36 + 18 + 12 = 120$ seconds. All approaches receive proportional service while guaranteeing at least 10 seconds to clear waiting queues.
\end{interviewq}

% -----------------------------------------------------------------------------
\section{Exercise 5: Classifying Congestion via Random Forest}
\label{sec:wb_congestion}
% -----------------------------------------------------------------------------

\begin{interviewq}{Whiteboard Problem 5: Congestion Classification Tracing}
\textbf{Problem Statement:}
In \texttt{backend/congestion\_predictor.py}, a Random Forest model classifies junction conditions into:
\texttt{LOW} (0), \texttt{MEDIUM} (1), \texttt{HIGH} (2), or \texttt{SEVERE} (3).
Given input vector: $\mathbf{x} = [\text{vehicle\_count} = 85, \text{avg\_speed} = 12.5 \text{ km/h}, \text{hour} = 18.5]$, explain:
\begin{enumerate}
    \item How the decision trees evaluate the vector.
    \item What classification is produced and why.
    \item How provenance metadata is attached in the response schema.
\end{enumerate}
\tcblower
\textbf{Step 1: Feature Evaluation Across Ensemble}
During peak evening rush hour ($18.5 = 6:30$ PM) with heavy vehicle presence (85 vehicles) and crawling speed (12.5 km/h):
\begin{itemize}
    \item Individual decision trees evaluate binary splits:
    \begin{align*}
    \text{Node 1:} &\quad (\text{vehicle\_count} > 60) \implies \text{True} \\
    \text{Node 2:} &\quad (\text{avg\_speed} \le 15.0) \implies \text{True} \\
    \text{Node 3:} &\quad (\text{hour} \ge 17 \land \text{hour} \le 20) \implies \text{True}
    \end{align*}
    \item Trees assign votes heavily biased toward Class 3 (\texttt{SEVERE}) and Class 2 (\texttt{HIGH}).
\end{itemize}

\textbf{Step 2: Majority Voting}
Suppose an ensemble of $B = 50$ trees produces:
\begin{itemize}
    \item Class 3 (\texttt{SEVERE}): 44 votes
    \item Class 2 (\texttt{HIGH}): 6 votes
    \item Class 1 (\texttt{MEDIUM}): 0 votes
    \item Class 0 (\texttt{LOW}): 0 votes
\end{itemize}
Argmax returns Class 3: \texttt{SEVERE}.

\textbf{Step 3: Response Schema and Provenance}
The backend formats the response matching \texttt{app.py} and \texttt{app\_demo.py}:
\begin{lstlisting}[language=json]
{
  "congestion_level": "SEVERE",
  "vehicle_count": 85,
  "average_speed_kmh": 12.5,
  "provenance": "synthetic_random_forest_model",
  "disclaimer": "Demonstration classifier trained on synthetic rule distributions."
}
\end{lstlisting}
The provenance flag explicitly discloses to the frontend that this prediction originates from synthetic demo distributions.
\end{interviewq}

% -----------------------------------------------------------------------------
\section{Exercise 6: Calculating Geodesic Haversine Distance}
\label{sec:wb_haversine}
% -----------------------------------------------------------------------------

\begin{interviewq}{Whiteboard Problem 6: Haversine Calculation}
\textbf{Problem Statement:}
In \texttt{backend/app.py} (\texttt{/api/driver\_advice}), the system calculates distance between a commuter and a traffic corridor. Given:
\begin{itemize}
    \item Commuter location $P_1 = (\phi_1, \lambda_1) = (12.9352^\circ\text{ N}, 77.6245^\circ\text{ E})$ (Koramangala)
    \item Silk Board Junction $P_2 = (\phi_2, \lambda_2) = (12.9172^\circ\text{ N}, 77.6228^\circ\text{ E})$
    \item Earth's mean radius $R = 6371.0$ km
\end{itemize}
Derive the great-circle distance $d$ in kilometers using the Haversine equation.
\tcblower
\textbf{Step 1: Convert Degrees to Radians}
\begin{align*}
\Delta \phi &= (12.9172 - 12.9352) \times \frac{\pi}{180} = -0.0180 \times 0.0174533 \approx -0.00031416 \text{ rad} \\
\Delta \lambda &= (77.6228 - 77.6245) \times \frac{\pi}{180} = -0.0017 \times 0.0174533 \approx -0.00002967 \text{ rad} \\
\phi_1 &= 12.9352 \times \frac{\pi}{180} \approx 0.22576 \text{ rad}, \quad \phi_2 = 12.9172 \times \frac{\pi}{180} \approx 0.22545 \text{ rad}
\end{align*}

\textbf{Step 2: Evaluate Haversine Formula}
\begin{align*}
a &= \sin^2\left(\frac{\Delta \phi}{2}\right) + \cos(\phi_1)\cos(\phi_2)\sin^2\left(\frac{\Delta \lambda}{2}\right) \\
\frac{\Delta \phi}{2} &= -0.00015708 \implies \sin^2\left(\frac{\Delta \phi}{2}\right) \approx 2.467 \times 10^{-8} \\
\frac{\Delta \lambda}{2} &= -0.000014835 \implies \sin^2\left(\frac{\Delta \lambda}{2}\right) \approx 2.201 \times 10^{-10} \\
\cos(\phi_1)\cos(\phi_2) &= \cos(0.22576)\cos(0.22545) \approx 0.9746 \times 0.9747 \approx 0.9499 \\
a &= 2.467 \times 10^{-8} + 0.9499 \times (2.201 \times 10^{-10}) \approx 2.488 \times 10^{-8}
\end{align*}

\textbf{Step 3: Compute Central Angle and Distance}
\begin{align*}
c &= 2 \cdot \arctan2(\sqrt{a}, \sqrt{1 - a}) \approx 2 \cdot \sqrt{a} \approx 2 \times 0.00015773 = 0.00031547 \text{ rad} \\
d &= R \times c = 6371.0 \times 0.00031547 \approx \mathbf{2.01 \text{ km}}
\end{align*}
\textbf{Conclusion:} Straight-line geographic separation is approximately $2.01$ kilometers.
\end{interviewq}

% -----------------------------------------------------------------------------
\section{Exercise 7: Tracing an End-to-End API Cycle}
\label{sec:wb_tracing}
% -----------------------------------------------------------------------------

\begin{interviewq}{Whiteboard Problem 7: Concurrent API Request Lifecycle}
\textbf{Problem Statement:}
Trace the complete request and response lifecycle when a user navigates to the Operations Dashboard (\texttt{/}), focusing on how \texttt{Dashboard.jsx} coordinates data fetching across \texttt{/api/status}, \texttt{/api/signal}, and \texttt{/api/violations}.
\tcblower
\textbf{Step 1: Component Mount and Polling Setup}
In \texttt{frontend/src/pages/Dashboard.jsx}, the component mounts. An effect hook sets up an interval timer:
\begin{lstlisting}[language=JavaScript]
useEffect(() => {
  fetchDashboardData();
  const interval = setInterval(fetchDashboardData, 2000);
  return () => clearInterval(interval);
}, []);
\end{lstlisting}

\textbf{Step 2: Concurrent Fetch via Promise.all}
In \texttt{fetchDashboardData()}, the client avoids waterfall latency by issuing three parallel asynchronous HTTP GET requests to the configured \texttt{VITE\_API\_BASE\_URL}:
\begin{lstlisting}[language=JavaScript]
const [statusRes, signalRes, violationsRes] = await Promise.all([
  axios.get(`${API_BASE}/api/status`),
  axios.get(`${API_BASE}/api/signal`),
  axios.get(`${API_BASE}/api/violations`)
]);
\end{lstlisting}

\textbf{Step 3: Server Processing and Serialization}
At the Flask backend (Render or local), the WSGI server handles requests concurrently:
\begin{itemize}
    \item \texttt{/api/status}: Reads thread-safe camera state or simulation generator, serializing total vehicle count, active lane counts, and FPS.
    \item \texttt{/api/signal}: Inspects active phase, remaining green countdown, and emergency preemption flags.
    \item \texttt{/api/violations}: Fetches recent violation records array.
\end{itemize}

\textbf{Step 4: Atomic Zustand Store Update}
Upon resolution of all three promises, the response data payloads are passed to \texttt{useStore}:
\begin{lstlisting}[language=JavaScript]
setDashboardData({
  status: statusRes.data,
  signal: signalRes.data,
  violations: violationsRes.data,
  isOffline: false
});
\end{lstlisting}
React triggers a single batched re-render of child components (\texttt{SignalView}, \texttt{MetricCards}, \texttt{ViolationTable}), preventing UI layout jitter.
\end{interviewq}

% -----------------------------------------------------------------------------
\section{Exercise 8: Debugging Backend Cold Starts and Network Timeouts}
\label{sec:wb_debugging}
% -----------------------------------------------------------------------------

\begin{interviewq}{Whiteboard Problem 8: Handling Server Cold-Start Failures}
\textbf{Problem Statement:}
On Render free-tier hosting, services spin down after 15 minutes of inactivity. When a user opens the dashboard, the first API request may take up to 50 seconds or time out. Describe how the frontend handles this without crashing.
\tcblower
\textbf{Step 1: Axios Timeout Configuration}
In our API client wrapper, requests specify a bounded timeout (e.g., 10 seconds):
\begin{lstlisting}[language=JavaScript]
const apiClient = axios.create({
  baseURL: import.meta.env.VITE_API_BASE_URL,
  timeout: 10000
});
\end{lstlisting}

\textbf{Step 2: Catching Network and Timeout Exceptions}
In \texttt{useStore.js}, the fetching action wraps calls in \texttt{try/catch}:
\begin{lstlisting}[language=JavaScript]
try {
  const data = await apiClient.get('/api/status');
  set({ status: data.data, isBackendAvailable: true, isLoading: false });
} catch (error) {
  set({ isBackendAvailable: false, isLoading: false });
}
\end{lstlisting}

\textbf{Step 3: Graceful UI Degradation}
Instead of throwing unhandled runtime exceptions or attempting to read properties of \texttt{undefined}:
\begin{itemize}
    \item The dashboard displays a prominent banner: \texttt{"Connecting to backend services... Free-tier server spinning up (30-50s)."}.
    \item Data visualization charts render fallback empty states or cached historical values.
    \item Polling continues automatically every 2 seconds until the server responds with HTTP 200, at which point the application seamlessly transitions to live state.
\end{itemize}
\end{interviewq}
